\begin{abstract}
Let \(a_1,\ldots,a_N\in\C\), let \(r\ge0\) be real
translation, and put
\[
 S_N=\sum_{n\le N}a_n,\qquad
 L_{r,N}=\sum_{n\le N}\frac{a_n}{r+n},\qquad L_{r,0}=0.
\]
The logarithmic primitive on a translated interval has an exact
Tauberian inverse that differs from the unshifted one:
\[
 \boxed{\quad
 S_N=(r+N)L_{r,N}-\sum_{m<N}L_{r,m},\qquad
 \frac{S_N}{N}
 =
 \left(1+\frac rN\right)L_{r,N}
 -\frac1N\sum_{m<N}L_{r,m}.
 \quad}
\]
Thus ordinary cancellation is equivalent to vanishing of an
\emph{anchored backward Ces\`aro defect}.  The extra term
\((r/N)L_{r,N}\) is not an error that may be dropped after translating
a physical fiber.

We prove the exact terminal criterion: once the unshifted signed
defect is \(o(1)\), ordinary cancellation holds if and only if
\((r/N)L_{r,N}=o(1)\).  Thus \(r=0\) needs no terminal condition;
for \(0<r=o(N)\) the exact scale is \(L_{r,N}=o(N/r)\), with
boundedness only a convenient stronger hypothesis; for
\(r\asymp N\) one needs \(L_{r,N}=o(1)\); and for \(r\gg N\) one
needs \(L_{r,N}=o(N/r)\).
The constant sequence gives the exact obstruction:
\[
 L_{r,N}=\sum_{n\le N}\frac1{r+n}\sim\frac Nr=o(1)
 \quad(r/N\to\infty),
 \qquad S_N/N=1.
\]
Hence even an unnormalized \(o(1)\) logarithmic statement can be
useless on a short translated fiber.

For a physical packet
\[
 Z_X=\sum_{\theta\in\Theta_X}
 c_{\theta,X}S_{\theta,X}(N_{\theta,X})
 +E_{\mathrm{ret},X},
\]
we obtain the exact signed aggregate formula
\[
 Z_X=\sum_{\theta\in\Theta_X}
 c_{\theta,X}N_{\theta,X}
 \cT_{\rho_{\theta,X},N_{\theta,X}}(L_{\theta,X})
 +E_{\mathrm{ret},X}.
\]
This identifies the quantitative certificate that a logarithmic
theorem must return after every literal origin, prefix, outer
coefficient, mask, and content term has been restored.  The identities
and counterexamples are \(\Lzero\); attachment to a complete native
packet is an \(\Lone\) certificate.  No uniform estimate for the
growing fixed-\(h_0\) affine M\"obius family is proved, so there is no
new \(\Ltwo\) arithmetic saving, parity breakthrough, or prime-pair
consequence.
\end{abstract}
